\section{Open Questions}

Five truly-open questions arising from this replication (\emph{not} already answered in the paper), each with a concrete, free-endpoint-only follow-up plan:

\begin{enumerate}

\item \textbf{Correlated noise vs.\ the paper's linear-in-$p$ law.}
Do the paper's linear-in-$p$ and $1/\sqrt{N}$ scaling laws survive under
correlated (non-Markovian, spatially/temporally correlated) noise, or do
they only hold for independent local depolarizing noise?
\emph{Basis:} both this replication and the paper use uncorrelated local
Pauli channels; real superconducting hardware exhibits crosstalk and $1/f$
drift that violates the Markovian assumption (Wallman \& Emerson 2016;
Blume-Kohout et al.\ 2022).
\emph{Next steps:} implement a Pauli--Lindblad correlated-noise model in
Qiskit-Aer via \texttt{qiskit-ibm-runtime} noise-learning primitives (free);
rerun the same 3-point $p$-sweep with (a) uncorrelated depolarizing
(baseline, done), (b) Pauli--Lindblad correlated across adjacent qubit
pairs, (c) same with slow $T_2^*$ drift; fit slope-in-$p$ for each and
report deviation. Cross-check against IBM's publicly released Kyiv/Sherbrooke
device noise-learning snapshots.

\item \textbf{Modern error mitigation vs.\ the 2021-era unmitigated numbers.}
How much of the depolarizing-noise energy error (0.04--0.22~Ha) can be
recovered by ZNE, PEC, or PEA that were not mature at the time of the paper?
\emph{Basis:} the 2022--2025 literature (Kandala et al.; Kim et al.\ 2023
\emph{Nature}; van den Berg et al.\ 2023) demonstrates ZNE and PEC recovery
to chemical accuracy at 1q errors $10^{-3}$--$10^{-4}$; not tested here.
\emph{Next steps:} rerun $p{=}10^{-3}$ and $p{=}10^{-2}$ with Mitiq's ZNE
(linear + Richardson, folding factors $[1,3,5]$) and PEC via Aer's analytic
depolarizing decomposition; report mitigated $|\Delta E|$ vs.\ unmitigated.
Predicted: ZNE recovers $p{=}10^{-3}$ to $5\times10^{-3}$~Ha; PEC to
$10^{-3}$~Ha at 10--100$\times$ sampling overhead.

\item \textbf{Transferability to strongly-correlated molecules.}
Do the conclusions transfer to LiH, N$_2$ stretched-bond, or H$_4$ chains
where the ground state has multireference character and depth-1
hardware-efficient ansatze demonstrably fail even in the ideal-simulator
limit?
\emph{Basis:} H$_2$/STO-3G is a benign 4-qubit test case; the paper's
scaling laws may not survive the qualitative regime change to
multireference systems where barren plateaus dominate.
\emph{Next steps:} build LiH/STO-3G (10 qubits, 631 JW Paulis) with the
same PySCF+OpenFermion pipeline; rerun the depolarizing sweep at
$p \in \{10^{-4},10^{-3}\}$; test whether $|\Delta E| \propto p$ still
holds. Bonus: stretched H$_4$ chain at $R{=}1.5$~\AA. Budget: $\sim$2~h
CPU on m1 for LiH; $\sim$10~h for H$_4$.

\item \textbf{Cross-noise-family transferability.}
How transferable are the scaling laws across noise-model families ---
depolarizing vs.\ amplitude-damping vs.\ coherent-rotation vs.\ realistic
$T_1/T_2$+gate-time IBMQ? The paper conflates these under ``noise'' but
they have qualitatively different failure modes (amplitude damping is
biased; coherent errors do not average to zero; depolarizing is unbiased).
\emph{Basis:} this replication tested only depolarizing; the paper tested
IBMQ device noise as a black box mixing all three.
\emph{Next steps:} split the sweep into (a) pure amplitude-damping
(\texttt{amplitude\_damping\_error}), (b) pure phase-damping, (c) coherent
over-rotation via a Kraus channel with small unitary error per gate;
compare slope-in-$p$ across families; check whether IBMQ device-noise
error decomposes as a weighted sum.

\item \textbf{ML-selected error mitigation across the paper's 12-ansatz family.}
Can a machine-learning-selected mitigation strategy (choose ZNE vs.\ PEC vs.\
symmetry-verification vs.\ Clifford-data-regression per (ansatz, noise,
molecule)) outperform any single mitigation, and does the selector
generalize across the paper's 12-ansatz family?
\emph{Basis:} the paper's C5 finding (``noise reshuffles ansatz ranking'')
plus modern QEM heterogeneity implies the optimal mitigation is
ansatz-dependent; no published work systematically learns this mapping.
\emph{Next steps:} first reproduce C5 (12-ansatz sweep at $p{=}10^{-3}$,
$\sim$1~h); apply ZNE, PEC, symmetry-verification (JW parity check) to
each; train a gradient-boosted classifier (features: depth, 2q-count,
entangling entropy at random init) to predict optimal mitigation;
cross-validate on LiH (question~3). Free endpoints only.

\end{enumerate}
